本研究では、ITOとspiro-OMeTAD薄膜の間に中間層を導入した場合の色変化の可逆性と酸化還元電流への影響を調べた。\\

色変化の可逆性については中間層を導入することで、spiro-OMeTADのみよりも向上するものは無かった。しかし、spiro-OMeTADのみとほとんど同じ可逆性を示すものもあり、中間層導入によるデメリットを最小限に抑えることができる可能性がある。\\
酸化還元電流については中間層を導入することで、spiro-OMeTADのみより大きくなった。\\
カフェ酸を導入すると、溶液濃度が濃くなるにつれて酸化電流は大きくなり、還元電流は小さくなった。2PACzを導入すると、溶液濃度が$\SI{30}{mmol/L}$まで酸化電流、還元電流どちらも大きくなった。\\
色変化の可逆性と酸化還元電流の比については、材料ごとに特徴が見られた。Br-2PACzとMe-2PACzではspiro-OMeTADのみに比べて酸化還元電流の比は向上したものの、色変化の可逆性は低下した。カフェ酸では、spiro-OMeTADのみに比べて酸化還元電流の比は低下したものの、色変化の可逆性は同等のものもあった。2PACzでは、spiro-OMeTADのみに比べて酸化還元電流の比は向上したものの、色変化の可逆性は一番いいものでも同等であった。\\
今後の展望としては、他の材料ではどのようになるのか、また酸化還元電流の大きさに双極子モーメントが関与しているかなど調べていく。